\exercisesection{}

\begin{enumerate}
\def\labelenumi{\arabic{enumi})}
\item
  设 G 为一个局部紧群，H 为 G 的闭子群。对 \(\xi\in H\)，置 \(\chi(\xi)=\Delta_{H}(\xi)\Delta_{G}(\xi)^{-1}\)。把 H 看作通过平移在 G 上右作用。证明在 \(\mathcal{K}^{\chi}(G)\) 上存在一个非零正线性型 I；除相差一个常数因子外，这样的线性型只有一个，而且它在左平移下不变。（使用命题 3，其中 X=G，并取 \(\mu\) 为 G 的左 Haar 测度。）
\item
  设 X 和 \(X'\) 是两个局部紧空间，局部紧群 H 在其中连续且适当地右作用。令 \(\theta\) 为从 X 到 \(X'\) 的适当连续映射，它与 H 的恒等映射相容（GT, III, §2, No.~4），并令 \(\theta': X/H \to X'/H\) 为由 \(\theta\) 通过取商得到的映射。设 \(f'\) 是 \(X'\) 上的连续函数，其支集与每个紧集的饱和集的交都是紧的。则 \(f = f' \circ \theta\) 在 X 中具有相同性质，并且
\end{enumerate}

\[
f ^ {b} = f ^ {\prime b} = \theta^ {\prime}
\]

（映射 \(f \mapsto f^b\)、\(f' \mapsto f'^b\) 均相对于在 H 上选定的同一个 Haar 测度。）

\begin{enumerate}
\def\labelenumi{\arabic{enumi})}
\setcounter{enumi}{2}
\tightlist
\item
  设 B 为局部紧空间，H 为局部紧群。置 \(\mathbf{X} = \mathbf{B} \times \mathbf{H}\)，群 H 通过下式在 X 上作用：
\end{enumerate}

\[
(b, \xi) \xi^ {\prime} = (b, \xi \xi^ {\prime}).
\]

令 \(\lambda\) 为 \(B = X / H\) 上的测度，\(\beta\) 为 \(H\) 上的左 Haar 测度。则 \(\lambda^{\sharp} = \lambda \otimes \beta\)。

\begin{enumerate}
\def\labelenumi{\arabic{enumi})}
\setcounter{enumi}{3}
\item
  设 X 是局部紧空间，局部紧群 H 在其中连续且适当地右作用。设 \(\beta\) 是 H 上的左 Haar 测度，\(\mu\) 是 X 上满足 \(\geqslant0\) 的测度。设 h 是 X 上满足 \(\geqslant0\) 的连续函数，且 \(h^{b}=1\)。若 f 是 X 上 \(\mu\)-可忽略的满足 \(\geqslant0\) 的函数，则函数 \((x,\xi)\mapsto f(x)h(x\xi)\) 在 X×H 上是 \((\mu\otimes\beta)\)-可忽略的。（存在递减的开集 \(\Omega_{1},\Omega_{2},\ldots\)，使得 \(\mu(\Omega_{n})\) 趋于 0，且 f 在 \(\Omega_{1}\cap\Omega_{2}\cap\cdots\) 之外为零。证明 \(\int^{*}\varphi_{\Omega_{n}}(x)h(x\xi)d\mu(x)d\beta(\xi)\leqslant\mu(\Omega_{n})\)，注意到函数 \((x,\xi)\mapsto\varphi_{\Omega_{n}}(x)h(x\xi)\) 是下半连续的。））
\item
  设 \(\mathbf{G}\) 为局部紧群。对每个 \(s \in \mathbf{G}\)，令 \(\psi_s\) 为加法群 \(\mathbf{R}\) 的自同构，由 \(\psi_s(x) = \Delta_{\mathbf{G}}(s)x\) 定义。令 \(\Gamma\) 为 \(\mathbf{G}\) 与 \(\mathbf{R}\) 的拓扑半直积，由 \(s \mapsto \psi_s\) 定义（GT, III, §2, No.~10）。证明 \(\Gamma\) 是幺模的。
\item
  设 \(G\) 为局部紧群，\(G_1, G_2, G_3\) 为闭子群，且 \(G_3 \subset G_1 \cap G_2\)。假设 \(G, G_1, G_2, G_3\) 都是幺模的，且 \(G_1 / G_3\) 具有有限总测度（对于每个在 \(G_1\) 下不变的测度）。令 \(\lambda, \mu, \nu\) 分别为 \(G / G_1, G / G_2, G_2 / G_3\) 上的不变测度，\(\varphi\) 为 \(G\) 到 \(G / G_1\) 的典范映射。设 \(f \in \mathcal{K}(G / G_1)\)。证明
\end{enumerate}

\[
a \int_ {\mathrm{G} / \mathrm{G} _ {1}} f (u) d \lambda (u) = \int_ {\mathrm{G} / \mathrm{G} _ {2}} d \mu (\dot {x}) \int_ {\mathrm{G} _ {2} / \mathrm{G} _ {3}} f (\varphi (x \xi)) d \nu (\dot {\xi})
\]

\((\dot{x} = x\mathrm{G}_{2}, \dot{\xi} = \xi \mathrm{G}_{3})\)，其中 \(a\) 是与 \(f\) 无关的常数。

¶ 7) a) 设 E 为局部紧空间，Γ 为在 E 上连续左作用的局部紧群。假设对每个 \(x \in E\)，从 Γ 到 E 的映射 \(s \mapsto s \cdot x\) 是适当的，并且 E 上存在一个在 Γ 下不变的非零有界正测度 μ。证明 Γ 是紧的。（设 \((s_n)\) 是 Γ 中点的一个序列。取 E 的一个不被 μ 忽略的紧子集 K。利用 Ch. IV, §4, Exer. 15，证明存在 \(x_0 \in E\)，使得对 n 的无穷多个值有 \(s_n x_0 \in K\)。~由此推出 \((s_n)\) 在 Γ 中有聚点。然后应用 GT, II, §4, Exer. 6。）

\begin{enumerate}
\def\labelenumi{\alph{enumi})}
\setcounter{enumi}{1}
\tightlist
\item
  设 G 为局部紧群，H 和 K 为 G 的两个闭子群。假设 G 和 H 是幺模的，且 G/H 对于 G 下的不变测度具有有限测度。证明 H 和 K 满足 GT, III, §4, Exer. 11 c) 的等价条件，当且仅当 K 是紧的。
\end{enumerate}

\begin{enumerate}
\def\labelenumi{\arabic{enumi})}
\setcounter{enumi}{7}
\tightlist
\item
  设 \(\mathfrak{G}\) 为局部紧群，G 和 H 为 \(\mathfrak{G}\) 的两个闭子群。假设每个 \(s \in \mathfrak{G}\) 都能且只能以一种方式写成 \(s = \xi x = y\eta\) \((\xi, \eta \in \mathbf{G}, x, y \in \mathbf{H})\)，并且 \(\xi, x, y, \eta\) 连续地依赖于 \(s\)。每个 \(\xi \in \mathbf{G}\) 定义 H 到 H 的同胚 \(\widehat{\xi}\)，其定义为 \(\xi x \in \widehat{\xi}(x)\mathbf{G} (x \in \mathbf{H})\)。每个 \(x \in \mathbf{H}\) 定义 G 到 G 的同胚 \(\widehat{x}\)，其定义为 \(\xi x \in \mathrm{H}\widehat{x}(\xi) (\xi \in \mathbf{G})\)。令 \(\mu, \alpha, \beta\) 分别为 \(\mathfrak{G}, \mathbf{G}, \mathbf{H}\) 上的左 Haar 测度。证明
\end{enumerate}

\[
\begin{array}{c} d \widehat {x} ^ {- 1} (\alpha) (\xi) = \frac {\Delta_ {\mathfrak {G}} (\widehat {\xi} (x)) \Delta_ {\mathrm{G}} (\widehat {x} (\xi))}{\Delta_ {\mathrm{H}} (\widehat {\xi} (x)) \Delta_ {\mathrm{G}} (\xi)} d \alpha (\xi), \\ d \widehat {\xi} ^ {- 1} (\beta) (x) = \frac {\Delta_ {\mathrm{G}} (\widehat {x} (\xi))}{\Delta_ {\mathfrak {G}} (\widehat {x} (\xi))} d \beta (x). \end{array}
\]

（取 \(f \in \mathcal{K}(G)\)。说明如下事实：借助 (23) 计算的 \(\int f(u) d\mu(u)\)，无论取 \(X = G\)、\(Y = H\)，还是取 \(X = H\)、\(Y = G\)，当 \(f\) 受到 \(G\) 中元素或 \(H\) 中元素的左平移时都保持不变。）

\begin{enumerate}
\def\labelenumi{\arabic{enumi})}
\setcounter{enumi}{8}
\tightlist
\item
  设 X 为局部紧空间，H 为在 X 上连续（因而适当地）右作用的紧群。记 \(\beta\) 为 H 上的归一化 Haar 测度，记 \(\pi: X \to X/H\) 为典范映射。令 \(\lambda\) 为 X/H 上的正测度，\(\lambda^{\sharp}\) 为 X 上对应的测度。
\end{enumerate}

\begin{enumerate}
\def\labelenumi{\alph{enumi})}
\item
  证明若 N 是 X/H 的 \(\lambda\)-可忽略子集，则 \(\overline{\pi}^{1}(N)\) 是 \(\lambda^{\sharp}\)-可忽略的。（计算集合 \(\overline{\pi}^{1}(U)\) 的测度，其中 U 是 X/H 中的开集。）
\item
  设 \(p\) 是一个有限数且 \(\geqslant 1\)，设 \(f\) 是 \(\mathcal{L}_{\mathrm{F}}^{p}(\mathrm{X},\lambda^{\sharp})\) 中的函数（F 为 Banach 空间，或 \(\mathrm{F} = \overline{\mathbf{R}}\)）。证明那些满足 \(x \in \mathrm{X}\) 且映射 \(\xi \mapsto f(x\xi)\) 关于 \(\beta\) 在 \(\mathrm{H}\) 上不可积的 x 所成的集合具有 \(\overline{\pi}^{1}(\mathrm{N})\) 的形式，其中 \(\mathrm{N}\) 是 \(\lambda\)-可忽略的。此外，若 \(f^b\) 是 \(\mathrm{X / H}\) 上的函数，在几乎处处（相对于 \(\lambda\)）定义，并满足 \(f^b (\pi (x)) = \int_{\mathrm{H}} f(x\xi) d\beta (\xi)\)，则 \(f^b\) 属于 \(\mathcal{L}_{\mathrm{F}}^{p}(\mathrm{X / H},\lambda)\)，且 \(\mathrm{N}_p(f^b)\leqslant \mathrm{N}_p(f)\)。
\item
  反之，若 \(g \in \mathcal{L}_{\mathrm{F}}^{p}(\mathrm{X / H}, \lambda)\)，则函数 \(g \circ \pi\) 属于 \(\mathcal{L}_{\mathrm{F}}^{p}(\mathrm{X}, \lambda^{\sharp})\)。若 \(p, q\) 是共轭指数，\(\mathrm{F}'\) 是 \(\mathrm{F}\) 的对偶，\(f \in \mathcal{L}_{\mathrm{F}}^{p}(\mathrm{X}, \lambda^{\sharp})\) 且 \(g \in \mathcal{L}_{\mathrm{F}'}^{q}(\mathrm{X / H}, \lambda)\)，则
\end{enumerate}

\[
\int_ {\mathrm{X}} \left\langle f (x), g (\pi (x)) \right\rangle d \lambda^ {\sharp} (x) = \int_ {\mathrm{X/H}} \left\langle f ^ {\flat} (z), g (z) \right\rangle d \lambda (z).
\]

\begin{enumerate}
\def\labelenumi{\arabic{enumi})}
\setcounter{enumi}{9}
\tightlist
\item
  采用 §1, No.~6, 命题 6 的记号；对每个 \(\alpha \in A\)，令 \(\mu_{\alpha}\) 为 \(G_{\alpha}\) 上的左 Haar 测度，从而 \(\mu_{\alpha}\) 的逆极限是左 Haar 测度 \(\mu\)，定义在 \(G\) 上。对每个 \(\alpha \in A\)，记 \(\lambda_{\alpha}\) 为 \(K_{\alpha}\) 上的归一化 Haar 测度。
\end{enumerate}

\begin{enumerate}
\def\labelenumi{\alph{enumi})}
\tightlist
\item
  设 \(p\) 是一个有限数且 \(\geqslant 1\)，设 \(f\) 是 \(\mathcal{L}_{\mathrm{F}}^{p}(\mathrm{G},\mu)\) 中的函数（F 为 Banach 空间，或 \(\mathrm{F} = \overline{\mathbf{R}}\)）。对每个 \(\alpha \in \mathbf{A}\)，函数
\end{enumerate}

\[
f _ {\alpha} (s) = \int_ {K _ {\alpha}} f (s \xi) d \lambda_ {\alpha} (\xi),
\]

在几乎处处（相对于 \(\mu\)）定义，属于 \(\mathcal{L}_{\mathrm{F}}^{p}(\mathrm{G},\mu)\)（练习 9）。证明相对于有向集 A，\(f_{\alpha}\) 在 \(p\) 阶均值意义下趋于 \(f\)（使用 §1, No.~6 的引理 2）。

\begin{enumerate}
\def\labelenumi{\alph{enumi})}
\setcounter{enumi}{1}
\tightlist
\item
  假设 A = N 且 p = 1。证明此时 \(f_{n}\) 在 G 中几乎处处（相对于 \(\mu\)）趋于 f（使用类似 Ch. V, §8, Exer. 16 中的方法）。
\end{enumerate}

¶ 11) 设 G 为局部紧群，H 为 G 的闭子群，λ 为 G 上的左 Haar 测度；假设 G/H 上存在一个在 G 下不变的测度 μ，使得 λ = μ \(^{\sharp}\) 且 μ(G/H) \textless{} +∞。令 ν 为 G/H 上满足 \(\nu (\mathrm{G / H}) < +\infty\) 的非零正测度。最后，令 \(h\) 为 \(\mathbf{G}\) 上具有 No.~4, 命题 8 所述性质的函数。

\begin{enumerate}
\def\labelenumi{\alph{enumi})}
\tightlist
\item
  设 A 为 G/H 的 Borel 子集。证明
\end{enumerate}

\[
\int_ {\mathrm{G}} h (s) \nu (s ^ {- 1} \mathrm{A}) d \lambda (s) = \nu (\mathrm{G/H}) \mu (\mathrm{A})
\]

（使用 No.~4 的命题 9 a)。）

\begin{enumerate}
\def\labelenumi{\alph{enumi})}
\setcounter{enumi}{1}
\tightlist
\item
  设 \(A_i\)（\(1 \leqslant i \leqslant n\)）为 \(G / H\) 的 Borel 子集，并对每个 \(i\) 令 \(b_i = \sup_{s \in G} \nu(s^{-1} A_i)\)。证明存在两个不同的指标 \(i, j\)，使得
\end{enumerate}

\[
\int_ {\mathrm{G}} h (s) \nu (s ^ {- 1} \mathrm{A} _ {i}) \nu (s ^ {- 1} \mathrm{A} _ {j}) d \lambda (s) \geqslant \frac {\left(\nu (\mathrm{G} / \mathrm{H})\right) ^ {2}}{\mu (\mathrm{G} / \mathrm{H})} \left(\left(\sum_ {k = 1} ^ {n} \mu (\mathrm{A} _ {k})\right) ^ {2} - \frac {\mu (\mathrm{G} / \mathrm{H})}{\nu (\mathrm{G} / \mathrm{H})} \sum_ {k = 1} ^ {n} b _ {k} \mu (\mathrm{A} _ {k})\right)
\]

（仿照 §1, Exer. 28 a) 论证。）

¶ 12) a) 设 X 为拓扑空间，\(f: X \to N\) 为上半连续函数。令 \(X_0\) 为 X 中 \(f\) 局部常值的点所成的集合，并令 \(Y_0 = X - X_0\)。按如下方式递归地定义 X 的两个子集序列 \((X_i)_{i \geqslant 0}\)、\((Y_i)_{i \geqslant 0}\)：\(X_i\) 是 \(Y_{i-1}\) 中使得 \(f|_{Y_{i-1}}\) 局部常值的点所成的集合，而 \(Y_i = Y_{i-1} - X_i\)。证明 \(X_i\) 是 \(Y_{i-1}\) 的稠密开子集，并且 \(\bigcap_i Y_i = \varnothing\)。

（证明 \(f(x) > i\) 对 \(Y_i\) 的每一点成立。）

\begin{enumerate}
\def\labelenumi{\alph{enumi})}
\setcounter{enumi}{1}
\tightlist
\item
  设 X 为集合，H 为在 X 上右作用的群，\(\pi\) 为 X 到 X/H 的典范映射，并令 A（分别地，B）为所有满足 \(\pi|A:A\to X/H\) 为单射（分别地，为满射）的 A ⊂ X 所成的集合。令 \((V_{1},V_{2},\ldots)\) 为由 A 中元素构成的 X 的可数覆盖。令
\end{enumerate}

\[
\mathrm{V} _ {i} ^ {\prime} = \mathrm{V} _ {i} \cap \complement \left(\mathrm{V} _ {1} \mathrm{H} \cup \mathrm{V} _ {2} \mathrm{H} \cup \dots \cup \mathrm{V} _ {i - 1} \mathrm{H}\right).
\]

证明 \(\mathrm{F} = \bigcup_{i\geqslant 0}\mathrm{V}_i^\prime \in \mathfrak{A}\cap \mathfrak{B}\)。

\begin{enumerate}
\def\labelenumi{\alph{enumi})}
\setcounter{enumi}{2}
\item
  设 X 为局部紧空间，H 为在 X 上连续且适当右作用的可数离散群，\(\pi\) 为 X 到 X/H 的典范映射，\(\mu\) 为 X 上满足 \(\geqslant 0\) 的测度。证明若 b) 中的集合 \(V_{i}\) 是可求积的（Ch. IV, §5, Exer. 17 d)），则 b) 中的集合 F 是一个可求积基本域。
\item
  维持 \(c\) ) 的假设，并使用 No.~10 的记号 \(\mathbf{H}_x, n(x)\)。若对某个 \(x \in X\) 存在 \(V\)，它是 \(x\) 在 \(X\) 中的邻域，使得 \(\mathrm{H}_y = \mathrm{H}_x\) 对所有 \(y \in V\) 成立，则称 x 为一般点。证明 \(X\) 的一般点正是函数 \(n\) 局部常值的点。证明一般点具有一个属于 \(\mathfrak{A}\) 且可求积的开邻域（使用 Ch. IV, §5, Exer. 17 d)）。
\item
  维持 \(c\) ) 的假设，并进一步假设 \(X\) 在无穷远处可数。证明若 \(X - X_0\) 是 \(\mu\)-可忽略的，则存在一个 Borel 且可求积的基本域 \(F\)。（将 \(a\) ) 的构造应用于函数 \(n\)。然后应用 \(c\) ) 和 \(d\) ) 的结果于 \(X_0\)。在每个 \(X_i\) 中作类似论证。）
\item
  设 \(U \in B\) 为开集。证明可以使 e) 中的集合 F 具有下列附加性质：1) \(F \subset U\)；2) 对 X 的每个紧子集 K，使得 Fs 与 K 相交的 \(s \in H\) 所成的集合是有限的。
\end{enumerate}

\begin{enumerate}
\def\labelenumi{\arabic{enumi})}
\setcounter{enumi}{12}
\tightlist
\item
  设 G 为紧群，\(\mu\) 为 G 上的 Haar 测度，u 为 G 的一个自同态，满足 \(u(\mathrm{G})\) 是 G 的开子群，且核 \(\overline{u}(e)\)（记为 \(G_{u}\)）是 G 的有限子群。
\end{enumerate}

\begin{enumerate}
\def\labelenumi{\alph{enumi})}
\item
  证明存在实数 \(h(u) > 0\) 及开邻域 \(U\)，它是 \(e\) 在 \(G\) 中的开邻域，使得对每个开集 \(V \subset U\)，\(u(V)\) 是 \(G\) 中的开集且 \(\mu(u(V)) = h(u)\mu(V)\)（参见 §1, No.~7, 命题 9 的推论）。
\item
  证明 \(h(u) = \operatorname{Card}(\mathrm{G} / u(\mathrm{G})) / \operatorname{Card}(\mathrm{G}_u)\)（使用 a) 和 No.~7 的命题 10，以两种方式计算 \(\mu(u(\mathrm{G}))\)）。
\end{enumerate}

